\section{Strong connection and principal comodule algebra}\label{sc}

Let $H$ be a Hopf algebra with coproduct~$\Delta$, counit~$\varepsilon$ and bijective
antipode~$S$. Next, let \mbox{$\delta\colon P\to P\otimes H$} be a coaction making $P$ a right
\mbox{$H$-comodule} algebra.
We shall frequently use
the Heyneman--Sweedler notation (with the summation sign suppressed) for coproduct and coaction:
\begin{gather*}
\Delta (h)=:h_{(1)}\otimes h_{(2)} ,\qquad
\delta(p)=:p_{(0)}\otimes p_{(1)} .
\end{gather*}
\begin{Definition}[\cite{hkmz11}]
Let $P$ be a right comodule algebra over a Hopf algebra
$H$ with bijective antipode, and let
\begin{gather*}
B:=P^{\operatorname{co}H} := \{ p \in P \, | \, \delta (p) = p \otimes 1 \}.
\end{gather*}
be the coaction-invariant
subalgebra. The comodule algebra $P$
is called {\em principal} if the following conditions are satisf\/ied:
\begin{enumerate}\itemsep=0pt
\item[1)]
the coaction of $H$ is Hopf--Galois, that is, the map
\begin{gather*}
\operatorname{can}_P \colon \  P \underset{B}{\otimes} P \longrightarrow P \otimes H,\qquad
	p \otimes q \longmapsto pq_{(0)} \otimes q_{(1)} ,
\end{gather*}
(called the canonical map) is bijective,
\item[2)]
the comodule algebra $P$ is right
$H$-equivariantly projective as a left  $B$-module, i.e.,
 there exists
 a right $H$-colinear and left $B$-linear
splitting of the multiplication map
$B \otimes P \rightarrow P$.
\end{enumerate}
\end{Definition}

\begin{Definition}[\cite{bh04}]\label{strong}
Let $ H$ be a Hopf algebra with  bijective antipode.
A \emph{strong connection $\ell$} on a right $H$-comodule algebra $P$ is a unital linear map $\ell\colon  H
\rightarrow  P \otimes  P$ satisfying
the following bicolinearity (i.e., left and right colinearity) and splitting conditions:
\begin{enumerate}\itemsep=0pt
\item[1)]
$(\operatorname{id}\otimes \delta) \circ
\ell = (\ell \otimes \operatorname{id}) \circ \Delta$,
$(\delta^L \otimes \operatorname{id}) \circ
\ell = (\operatorname{id} \otimes \ell) \circ
\Delta$, where
$\delta^L :=(S^{-1}\otimes\operatorname{id})\circ\mathrm{f\/lip}\circ\delta$;
\item[2)]
$\widetilde{\operatorname{can}_P} \circ \ell=1 \otimes \operatorname{id}$, where
$\widetilde{\operatorname{can}_P}\colon  P\otimes  P\ni p\otimes q\mapsto (p\otimes 1)\delta(q)\in  P\otimes H$.
\end{enumerate}
\end{Definition}

We will use the  Heyneman--Sweedler-type notation
$
\ell(h)=:\ell(h)^{\langle 1 \rangle}\otimes
 \ell(h)^{\langle 2 \rangle}
$
 with the summation sign suppressed.
 One can easily prove (see \cite[p.~599]{hkmz11} and references therein) that
a~comodule algebra is principal if and only if it admits a strong
connection.
Here we need the following slight generalization of this fact:
\begin{Lemma}\label{general}
Let $H$ be a Hopf algebra with bijective antipode. Then a right $H$-comodule algebra~$P$ is principal if and
only if it admits a~$($not necessarily unital$)$ linear map $\ell \colon H \rightarrow  P \otimes  P$ satisfying the
bicolinearity and splitting conditions of Definition~{\rm \ref{strong}}.
\end{Lemma}

\begin{proof}
Assume that $\ell \colon H \rightarrow  P \otimes  P$ is a linear map satisfying the bicolinearity and splitting
conditions.
Let $\pi_B\colon P\otimes P\to P\otimes_BP$ be the canonical surjection. Def\/ine
\begin{gather*}
L\colon \ P\otimes H\ni p\otimes h\longmapsto \pi_B\big(p \ell(h)^{\langle 1\rangle}\otimes\ell(h)^{\langle 2\rangle}\big)\in P\underset{B}{\otimes}P.
\end{gather*}
It follows immediately from the  splitting property of $\ell$ that
\begin{gather*}
\operatorname{can}_P(L(p\otimes h))= p \, \widetilde{\operatorname{can}_P}(\ell(h))=p\otimes h.
\end{gather*}
Applying $\operatorname{id}\otimes\varepsilon$ to the splitting condition for $\ell$, we obtain $m\circ\ell=\varepsilon$, where $m$ is
the multiplication map on~$P$.
Combining it with the left colinearity of $\ell$, which implies that
\begin{gather*}
\delta\big(q_{(0)}\ell(q_{(1)})^{\langle 1\rangle}\big)\otimes\ell(q_{(1)})^{\langle 2\rangle}
=q_{(0)}\ell(q_{(3)})^{\langle 1\rangle}\otimes q_{(1)}S(q_{(2)})\otimes\ell(q_{(3)})^{\langle 2\rangle}\\
\hphantom{\delta\big(q_{(0)}\ell(q_{(1)})^{\langle 1\rangle}\big)\otimes\ell(q_{(1)})^{\langle 2\rangle}}{}
=q_{(0)}\ell(q_{(1)})^{\langle 1\rangle}\otimes 1\otimes\ell(q_{(1)})^{\langle 2\rangle},
\end{gather*}
we obtain
\begin{gather*}
L(\operatorname{can}_P(p\otimes q)) = \pi_B \big(pq_{(0)} \ell(q_{(1)})^{\langle 1\rangle}
 \otimes  \ell(q_{(1)})^{\langle 2\rangle}\big)
 = \pi_B \big(p \otimes q_{(0)} \ell(q_{(1)})^{\langle 1\rangle} \ell(q_{(1)})^{\langle 2\rangle}\big)
 \\
\hphantom{L(\operatorname{can}_P(p\otimes q))}{}
=  \pi_B(p\otimes q).
\end{gather*}
Finally, the bicolinearity of $\ell$ implies that the formula $s(p):= p_{(0)}\ell(p_{(1)})$ def\/ines a left $B$-linear
right $H$-colinear splitting
of the multiplication map $B\otimes P\to P$. The reverse implication (principality $\Rightarrow$ existence of a~bicolinear~$\ell$ with the splitting property)
follows from~\cite[Lemma~2.2]{bh04}.
\end{proof}


\section[Join and fusion of unital $C^*$-algebras]{Join and fusion of unital $\boldsymbol{C^*}$-algebras}


Let $\otimes_{\min}$ denote the spatial (minimal) tensor product.
Due to the fact that minimal tensor products preserve injections
 (e.g., see \cite[Proposition~4.22]{t-m79} or
\cite[Section~1.3]{w-s94}),
for any unital $C^*$-algebras $A_1$ and $A_2$ the natural maps
\begin{gather*}
A_1\ni a\longmapsto a\otimes 1\in A_1 \underset{\min}{\otimes} A_2\qquad \text{and}
\qquad  A_2\ni a\longmapsto 1\otimes a\in A_1 \underset{\min}{\otimes} A_2
\end{gather*}
are injective. This allows us to view $A_1$ and $A_2$ as subalgebras of
$A_1 \otimes_{\min} A_2$. We use this natural identif\/ication in what follows.
\begin{Definition}\label{gjoin}
Let $A_1$ and $A_2$  be unital $C^*$-algebras, and
let $C$ be a unital $C^*$-algebra equipped with two $C^*$-ideals $J_1$ and $ J_2$
such that the direct sum
$
\pi_1\oplus\pi_2\colon C\rightarrow (C/J_1) \oplus (C/J_2)
$
of the canonical surjections $\pi_i\colon C\to C/J_i$, $i\in \{1,2\}$, is surjective.
We call the unital $C^*$-algebra
\begin{gather*}
A_1 \underset{\pi}{\circledast} A_2 :=
\big\{x\in C \underset{\min}{\otimes} A_1 \underset{\min}{\otimes} A_2 \,
\big|\,
(\pi_i\otimes \operatorname{id})(x) \in
(C/J_i)\otimes_{\min}A_i
, \, i\in\{1,2\}
\big\}
\end{gather*}
the \emph{fusion} $C^*$-algebra of $A_1$ and~$A_2$ over~$C$.
When $C$ is the $C^*$-algebra $C([0,1])$ of all continuous complex-valued
functions on the unit interval and
$\pi_1:=\mathrm{ev}_1$ and $\pi_2:=\mathrm{ev}_0$ are
the evaluation maps respectively at $1$ and $0$,
we call $A_1 \circledast_\pi A_2$ the \emph{join} $C^*$-algebra of $A_1$ and~$A_2\,$, and denote it by
 $A_1 \circledast A_2$.
\end{Definition}
